\chapter{การประเมินผล}

ในบทนี้จะกล่าวถึงกระบวนการประเมินผลและการทดสอบเกมต้นแบบที่พัฒนาขึ้น เพื่อวิเคราะห์ว่าเกมมีคุณสมบัติสอดคล้องกับวัตถุประสงค์และแนวทางการออกแบบที่กำหนดไว้หรือไม่ โดยการประเมินผลจะพิจารณาจากการทดสอบการเล่นจริง (Playtest) และการทดสอบโดยผู้ใช้งาน (User Testing) รวมถึงการรวบรวมความคิดเห็นและข้อเสนอแนะจากผู้ทดสอบ เพื่อนำไปใช้ในการปรับปรุงและพัฒนาเกมให้มีประสิทธิภาพมากยิ่งขึ้น

\section{วิธีการกำหนดเกณฑ์การประเมิน}

การประเมินผลของโครงงานนี้กำหนดเกณฑ์ในการประเมินจากองค์ประกอบสำคัญของเกม ได้แก่ ความสามารถในการทำงานของระบบเกม ความสนุกของรูปแบบการเล่น และความสะดวกในการควบคุมตัวละคร โดยมีเกณฑ์การประเมินดังต่อไปนี้

\begin{itemize}
\item ความถูกต้องของระบบเกม เช่น ระบบการเคลื่อนที่ ระบบการต่อสู้ และระบบศัตรูสามารถทำงานได้ตามที่ออกแบบไว้
\item ความเสถียรของโปรแกรม เช่น เกมสามารถทำงานได้ต่อเนื่องโดยไม่เกิดข้อผิดพลาดร้ายแรง
\item ความเข้าใจง่ายของการควบคุมเกม เช่น ผู้เล่นสามารถเรียนรู้วิธีการควบคุมได้ง่าย
\item ความสนุกและความท้าทายของเกม เช่น รูปแบบการเล่นมีความน่าสนใจและมีความท้าทายที่เหมาะสม
\end{itemize}

เกณฑ์ดังกล่าวถูกนำมาใช้ในการประเมินผลการทดสอบเกมต้นแบบจากผู้เล่นจริง

\section{วิธีการทดสอบ (Playtest และ User Testing)}

การทดสอบเกมในโครงงานนี้แบ่งออกเป็นสองรูปแบบหลัก ได้แก่ การทดสอบโดยผู้พัฒนา (Playtest) และการทดสอบโดยผู้ใช้งาน (User Testing)

การทดสอบแบบ Playtest เป็นการทดสอบโดยทีมผู้พัฒนา เพื่อทดสอบการทำงานของระบบต่าง ๆ ภายในเกม เช่น ระบบการเคลื่อนที่ของตัวละคร ระบบการโจมตี ระบบศัตรู และระบบไอเทม เพื่อค้นหาข้อผิดพลาดของระบบและปรับปรุงให้ระบบสามารถทำงานได้อย่างถูกต้อง

การทดสอบแบบ User Testing เป็นการให้ผู้เล่นทดลองเล่นเกมต้นแบบ เพื่อประเมินประสบการณ์การเล่นของผู้ใช้ โดยผู้ทดสอบจะทดลองเล่นเกมในด่านแรกของเกม และให้ข้อเสนอแนะเกี่ยวกับความสนุก ความยากง่ายของเกม และความเข้าใจในการควบคุมตัวละคร

\section{ผลการทดสอบและข้อเสนอแนะ}

จากการทดสอบเกมต้นแบบ พบว่าระบบหลักของเกมสามารถทำงานได้ตามที่ออกแบบไว้ ผู้เล่นสามารถควบคุมตัวละคร เคลื่อนที่ ต่อสู้กับศัตรู และเล่นมินิเกมบนเวทีได้ตามรูปแบบการเล่นที่กำหนดไว้

อย่างไรก็ตาม จากการทดสอบโดยผู้ใช้งานพบข้อเสนอแนะบางประการ เช่น การปรับสมดุลของความยากของศัตรูในบางช่วงของเกม การปรับปรุงการตอบสนองของการควบคุมตัวละครให้มีความลื่นไหลมากขึ้น และการเพิ่มคำแนะนำภายในเกมเพื่อช่วยให้ผู้เล่นเข้าใจระบบการเล่นได้ง่ายขึ้น

ข้อเสนอแนะดังกล่าวถูกนำไปใช้เป็นแนวทางในการปรับปรุงระบบเกมในขั้นตอนการพัฒนาต่อไป

\section{สรุปผลการประเมินและแนวทางพัฒนาเพิ่มเติม}

จากผลการทดสอบและการประเมินเกมต้นแบบ สามารถสรุปได้ว่าเกมที่พัฒนาขึ้นสามารถทำงานได้ตามวัตถุประสงค์ของโครงงาน โดยระบบหลักของเกม เช่น ระบบการต่อสู้ ระบบศัตรู และระบบการสำรวจ สามารถทำงานได้อย่างถูกต้อง

อย่างไรก็ตาม ยังมีแนวทางในการพัฒนาเพิ่มเติมในอนาคต เช่น การเพิ่มเนื้อหาของเกมให้มีความหลากหลายมากขึ้น การปรับปรุงสมดุลของเกม และการพัฒนาระบบกราฟิกและเสียงให้มีความสมบูรณ์มากยิ่งขึ้น เพื่อให้เกมมีประสบการณ์การเล่นที่ดียิ่งขึ้นสำหรับผู้เล่น